\subsection{導入：プロセスを学ぶ理由}
ソフトウェア工学プロセスは、開発手順の名前を選ぶだけの問題ではない。要求を理解し、設計し、実装し、検証し、移行し、運用し、保守し、廃止するまでには、多数の活動が相互に関係する。ある活動の出力は、別の活動の入力になる。したがって、単独の作業だけでなく、作業同士の接続を設計する必要がある。

SWEBOK 第10章は、プロセスを「概念」「ライフサイクル」「プロセス評価と改善」の三つの視点から整理する。ソフトウェア工学の世界では、プロセス標準化に関する議論が活発であり、ISO/IEC/IEEE 12207 のような標準も重要な参照点になる。

重要なのは、理想的な唯一のプロセスは存在しないという点である。プロセスは、プロジェクトの性質、製品の規模、利害関係者、組織の成熟度、変更の頻度、求められる品質水準に合わせて選択・適応・運用される。さらに、その運用は経験的測定によって支えられるべきである。

\begin{pointbox}{この章の読み方}本資料では、専門用語は原則として日本語で示し、必要に応じて英語を括弧内に併記する。英語名は、元資料やツール上の名称を確認するための手がかりである。説明は初学者向けに再構成しているため、原文の順序を保ちつつ、実務での見方を補足している。\end{pointbox}
